\chapter{合并全部范围，完成定理}\label{sec:completion}

\section{现在已经建立的依赖关系}
我们已经分别完成必要条件、阈值下界和充分性的所有参数段。
下面将它们合并。合并时应按范围使用结论，不能从相邻位置数的可解性自动推出中间维数。
这里没有增加或删除位置的默认规则；每个被引用的构造都覆盖其指定的原位置集合。

\begin{center}
\begin{tabular}{p{.36\textwidth}p{.48\textwidth}}
\toprule 证明任务 & 已建立的结论\\\midrule
任意可解输入的必要条件 & 第\ref{sec:obstruction}讲的共同余数障碍\\
最优阈值的下界 & 第\ref{sec:obstruction}讲的 $2p$ 元反例\\
宽维数与独立含零终端 & 第\ref{ch:descent-finish}、\ref{sec:blocks}、\ref{sec:wide}讲\\
一般输入到核心 & 第\ref{sec:compression}讲\\
核心内的共同算术终端 & 第\ref{sec:arithmetic}--\ref{sec:controllers}讲\\
各参数的实际矩阵资源 & 第\ref{ch:odd}--\ref{ch:upper}讲\\
三个最小奇素数的补充 & 第\ref{ch:three}、\ref{app:small}讲\\
\bottomrule
\end{tabular}
\end{center}

\section{主定理的完整证明}

\begin{proof}[定理\ref{thm:prime}的证明]
必要性由引理\ref{lem:obstruction}给出，下界由命题\ref{prop:lower}给出。对于充分性，$p\ge11$ 时：$n\ge4p$ 用命题\ref{prop:wide}；$n=3p$ 用命题\ref{prop:threep}；其余临界带先用命题\ref{prop:entry}压缩，再按下表完成。
\begin{center}
\begin{tabular}{ll}
\toprule 参数范围 & 核内结论\\\midrule
$n=2p+r,\ 1\le r<p,\ r$ 奇 & 命题\ref{prop:odd}\\
$r=2$ & 命题\ref{prop:r2}\\
$4\le r<p,\ r$ 偶 & 命题\ref{prop:even}\\
$3p<n<4p$ & 命题\ref{prop:upper}\\\bottomrule
\end{tabular}
\end{center}
$p=5$ 时，$11\le n\le15$ 已由上述下半带及等块结论处理；$n=16,17,18$ 用引理\ref{lem:collision-small}；$n=19$ 满足引理\ref{lem:integer-range}；$n\ge20$ 用命题\ref{prop:wide}。
$p=7$ 时，下半带处理 $15\le n\le21$；$n=22$ 用引理\ref{lem:seven22}；$23\le n\le26$ 用引理\ref{lem:collision-small}；$n=27$ 满足引理\ref{lem:integer-range}；更大维数用命题\ref{prop:wide}。
$p=3$ 时，$n=7$ 用命题\ref{prop:entry}及\ref{prop:odd}，$n=8,10$ 用引理\ref{lem:eight}、\ref{lem:ten}，$n\ge11$ 用引理\ref{lem:triple-large}。$n=9$ 将位置排成 $3\times3$，先平均各行再平均各列，六次操作即达到全局均值。
这些范围覆盖全部奇素数及全部 $n\ge2p+1$。推论\ref{cor:threshold}随定理和第\ref{sec:obstruction}讲的二平均论证成立。
\end{proof}

\section{从原始有理输入使用这个判据}
给定 $x\in\Q^n$，先计算全局均值 $\bar x$。
若配置本来恒定，任务已经完成。否则减去 $\bar x$，
乘共同分母后除去坐标公因子，得到本原零和整数代表 $X$。
再计算全部坐标差的最大公因子 $G$。
当奇素数 $p$ 满足 $n\ge2p+1$ 时，$G$ 仅含 $p$ 因子正是可解条件。

证明同时解释了两个方向。若 $G$ 含其他素因子，该素数处的共同非零余数永远保留，
因而任何操作选择都不可能成功。
若 $G$ 只含 $p$ 因子，则根据 $(p,n)$ 落入的范围，
整数下降或核心构造给出一条有限平均词，最后回到原尺度，精确达到 $\bar x\mathbf1$。

整份证明的终点始终是实际位置上的有限序列。
二维矩阵、局部化环和有限余数类所做的工作，是把这条序列的存在组织成可以证明的步骤。
判据本身十分简短，其充分性背后的结构则由本讲义逐层展开。


